\section{Homogeneous replacement on a coordinate ball}
\label{sec:local-homogeneous-dirichlet}

To show that the Monge--Amp\`ere measure of an envelope vanishes off the contact set, we shall replace the envelope on small balls by a solution of the homogeneous equation with the same boundary values. We first recall the smooth Dirichlet problem with a smooth positive density; the replacement is obtained by letting a constant density tend to zero, and continuous boundary values are then handled by uniform approximation.

In this section we work on the ball $B=B_R=\{|z|<R\}\subset\mathbb C^n$, put $\beta=\ddc|z|^2$, and let $\theta$ be any smooth real $(1,1)$-form defined near $\overline B$. Distances and the Laplacian $\Delta$ are Euclidean.

We use smooth comparison on subdomains $D$ of the ball. Let $a,b$ be smooth functions on $\overline D$ with $\theta+\ddc a>0$ and $\theta+\ddc b>0$. If $(\theta+\ddc a)^n\geq(\theta+\ddc b)^n$ in $D$ and $a\leq b$ on $\partial D$, then $a\leq b$ in $D$. Indeed, integrating the derivative of the logarithm of the determinant along the segment between the two positive matrices gives a linear elliptic inequality for $a-b$, and the maximum principle proves the assertion. The boundary of $D$ need not be smooth for this argument.

\begin{lemma}\label{lem:ball-smooth-dirichlet}
Let $\varphi\in C^\infty(\partial B)$, and let $F$ be a smooth strictly positive function on $\overline B$. There is a unique smooth function $v$ on $\overline B$ such that
\[
 \theta+\ddc v>0,\qquad
 (\theta+\ddc v)^n=F\beta^n,
 \qquad v|_{\partial B}=\varphi.
\]
\end{lemma}
\begin{proof}
Extend $\varphi$ to a smooth function $\widetilde\varphi$ on $\overline B$. For sufficiently large $A$, the function $\underline v=\widetilde\varphi+A(|z|^2-R^2)$ is a smooth strict subsolution with the prescribed boundary values: $\theta+\ddc\underline v$ is positive and its determinant is at least $(\max_{\overline B}F)\beta^n$. For $n\geq2$, the classical smooth Dirichlet theorem of Guan--Li therefore applies on the ball, with the arbitrary smooth background $\chi=\theta$ and the smooth positive right-hand side $F$. Their operator is $(\mathrm{i}/2)\partial\bar\partial$, so our potential is $\pi$ times theirs. No closedness condition is needed here; such conditions enter only in their separate theorems on compact manifolds. For $n=1$ the lemma is the Dirichlet problem for the Poisson equation. Uniqueness follows from smooth comparison, applied in both directions.
\end{proof}

\Needspace{15\baselineskip}
\begin{theorem}\label{thm:ball-homogeneous-replacement}\index{Dirichlet problem!homogeneous equation}
For every $\varphi\in C(\partial B)$ there exists a continuous $\theta$-plurisubharmonic function $H_\varphi$ on $\overline B$ such that
\[
 (\theta+\ddc H_\varphi)^n=0\quad\hbox{on }B,
 \qquad H_\varphi|_{\partial B}=\varphi.
\]
It dominates every continuous $\theta$-plurisubharmonic function $p$ on $\overline B$ with $p|_{\partial B}\leq\varphi$. For $\varphi,\psi\in C(\partial B)$, the construction satisfies
\[
 \|H_\varphi-H_\psi\|_{C(\overline B)}
       \leq\|\varphi-\psi\|_{C(\partial B)}.
\]
For smooth boundary data the constructed function is Lipschitz on $\overline B$.
\end{theorem}
\begin{proof}
\emph{Step 1: boundary barriers.} We first take $\varphi$ smooth and $0<\varepsilon\leq1$, and let $v_\varepsilon$ be the solution of \cref{lem:ball-smooth-dirichlet} with the constant density $F\equiv\varepsilon$. The strict subsolution $\underline v$ constructed there for $F\equiv1$ lies below every $v_\varepsilon$, by smooth comparison. Choose a smooth function $q$ with $q=\varphi$ on the boundary and $\Delta q=-C$, where $C$ is so large that $\theta+\ddc v_\varepsilon\geq0$ implies $\Delta v_\varepsilon\geq-C$. The maximum principle for the Laplacian gives $v_\varepsilon\leq q$. Both barriers are smooth and equal $\widetilde\varphi$ on the boundary. Consequently
\begin{equation}\label{eq:ball-boundary-distance}
 |v_\varepsilon(z)-\widetilde\varphi(z)|
       \leq C_1\operatorname{dist}(z,\partial B),
\end{equation}
with $C_1$ independent of $\varepsilon$.

\emph{Step 2: a Lipschitz bound independent of $\varepsilon$.} Let $h\in\mathbb C^n$ be small, and set $D_h=\{z\in B:z+h\in B\}$. Since $\theta$ is smooth, there is $C_2>0$, independent of the translation $h$, such that
\[
 \theta(z)-\theta(z+h)+C_2|h|\beta\geq0.
\]
For a constant $C_3>0$ to be fixed, define on $D_h$
\[
 w(z)=v_\varepsilon(z+h)
          +C_2|h|(|z|^2-R^2)-C_3|h|.
\]
Then $\theta+\ddc w\geq\theta(z+h)+\ddc v_\varepsilon(z+h)>0$, and the determinant of $\theta+\ddc w$ is at least $\varepsilon\beta^n$, because the Euclidean volume form is translation invariant. On $\partial D_h$, one of the points $z$ and $z+h$ lies on $\partial B$, and the other has distance at most $|h|$ from $\partial B$. Estimate \eqref{eq:ball-boundary-distance} and the Lipschitz bound for $\widetilde\varphi$ therefore give $w\leq v_\varepsilon$ on $\partial D_h$ if $C_3$ is chosen sufficiently large, and smooth comparison extends this inequality to all of $D_h$. Together with the same argument for $-h$, this gives
\[
 |v_\varepsilon(z+h)-v_\varepsilon(z)|\leq C_4|h|,
\]
with constants independent of $\varepsilon$. The same bound for arbitrary pairs of points follows by subdividing the segment joining them inside the convex ball, and it extends to $\overline B$ because each $v_\varepsilon$ is continuous there.

\emph{Step 3: the limit $\varepsilon\to0$.} Smooth comparison also shows that $v_\varepsilon$ increases as $\varepsilon$ decreases. By the barriers and the common Lipschitz bound, the limit $H_\varphi$ is Lipschitz and the convergence is uniform on $\overline B$. The inequality $\theta+\ddc v_\varepsilon>0$ of currents passes to the limit, and continuity of local bounded products under uniform convergence gives $(\theta+\ddc H_\varphi)^n=0$. For smooth boundary values $\varphi,\psi$, denote the solutions of \cref{lem:ball-smooth-dirichlet} with $F\equiv\varepsilon$ by $v_{\varepsilon,\varphi}$ and $v_{\varepsilon,\psi}$. Smooth comparison gives $\|v_{\varepsilon,\varphi}-v_{\varepsilon,\psi}\|_\infty\leq\|\varphi-\psi\|_\infty$, and passing to the limit proves the contraction for smooth data.

\emph{Step 4: continuous boundary values.} For continuous $\varphi$, choose smooth $\varphi_j\to\varphi$ uniformly on $\partial B$. By the contraction, $H_{\varphi_j}$ is uniformly Cauchy. Its limit $H_\varphi$ is continuous and $\theta$-plurisubharmonic, has the prescribed boundary values, and satisfies the homogeneous equation, again by continuity of bounded products under uniform convergence. Passing to the limit also proves the contraction for all continuous boundary data.

\emph{Step 5: domination.} It remains to prove the domination property, which does not follow from a comparison of total masses. First take $\varphi$ smooth and $p$ as in the statement. For $\delta>0$ put $p_\delta=p+\delta(|z|^2-R^2)$, so that $\theta+\ddc p_\delta\geq\delta\beta$. By \cref{lem:qpsh-upper-test}, applied with $\theta-\delta\beta$ in place of $\theta$, every $C^2$ function $a$ touching $p_\delta$ from above at an interior point $z_0$ satisfies
\[
 \theta(z_0)+\ddc a(z_0)\geq\delta\beta.
\]
Suppose that $p_\delta-v_\varepsilon$ had a positive maximum. Since this difference is nonpositive on the boundary, the maximum would be attained at an interior point, where $v_\varepsilon$ plus the maximum touches $p_\delta$ from above. The determinant of $\theta+\ddc v_\varepsilon$ at that point would then be at least $\delta^n\beta^n$, contradicting the equation for $v_\varepsilon$ when $\varepsilon<\delta^n$. Hence $p_\delta\leq v_\varepsilon$ for small $\varepsilon$. Letting first $\varepsilon$ and then $\delta$ decrease to zero, we obtain $p\leq H_\varphi$.

Finally, for continuous $\varphi$, take the smooth approximants $\varphi_j$ above and set $a_j=\|\varphi_j-\varphi\|_\infty$. The boundary values of $p-a_j$ are at most $\varphi_j$, so the result for smooth data gives $p-a_j\leq H_{\varphi_j}$. Uniform convergence proves $p\leq H_\varphi$.
\end{proof}
